% ============================================================
%  L-Shaped Outer-Approximation Algorithm for the HIV EID VRP
%  Paste this into an Overleaf article / section.
%  Required packages (add to preamble if not already present):
%    \usepackage{amsmath,amssymb,amsthm}
%    \usepackage{algorithm}
%    \usepackage{algpseudocode}
%    \usepackage{booktabs}
%    \usepackage{xcolor}
% ============================================================

\section{Solution Algorithm: Outer-Approximation L-Shaped Method}
\label{sec:algorithm}

\subsection{Problem Structure and Motivation}

The full objective for the Average (or Maximum) Turnaround-Time
minimisation problem contains a nonlinear term arising from the
laboratory queuing delay.  Specifically, the total expected waiting
time at the laboratory (W_{\text{lab}}) can be decomposed as

\begin{equation}
  W_{\text{lab}}
  \;=\;
  \underbrace{%
    \frac{K-1}{2K\mu\rho}
    + \frac{SCV \cdot \rho}{2\mu(1-\rho)}
    + \frac{1}{\mu}
  }_{\text{constant (given }\rho\text{)}}
  \;+\;
  \underbrace{%
    \frac{1}{2K^{2}\mu^{2}(1-\rho)}
    \sum_{r \in \mathcal{R}} \frac{D_r^{2}}{H_r}
  }_{\varphi:\;\text{route-dependent}},
  \label{eq:wlab}
\end{equation}

where $\rho = \sum_{i\in\mathcal{C}} \lambda_i / (K\mu)$ is the
laboratory utilisation (a constant once the clinic set is fixed),
$D_r = H_r \sum_{i \in \mathcal{C}_r} \lambda_i$ is the total demand
batched per vehicle visit on route $r$, and $H_r = L_r / z_r$ is the
headway (time between successive vehicle visits to any point on route
$r$).  The first three terms of \eqref{eq:wlab} are independent of the
routing decisions once $\rho$ is known; only $\varphi$ depends on how
clinics are assigned to routes and how many vehicles serve each route.

The term $D_r^{2}/H_r$ enters the mixed-integer programme as a
rotated second-order cone (SOC) constraint
\begin{equation}
  D_r^{2} \;\leq\; H_r \cdot \beta_r, \quad r \in \mathcal{R},
  \label{eq:soc}
\end{equation}
with $\beta_r$ appearing in the objective.  This makes the master
problem a Mixed-Integer Quadratically Constrained Programme (MIQCP),
which is substantially harder to solve than a pure Mixed-Integer Linear
Programme (MILP).  Furthermore, the nonconvex interaction between the
binary clinic-assignment variables $y_{ir}$, the integer vehicle-count
variables $z_r$, and the continuous headway $H_r$ leads to long
solution times that grow rapidly with the number of clinics.

\subsection{Outer-Approximation Reformulation}

We exploit the fact that $f(D,H) = D^{2}/H$ is jointly convex in
$(D,H)$ for $H > 0$.  Its epigraph can therefore be outer-approximated
by a (potentially infinite) family of linear supporting hyperplanes.

\begin{proposition}[Supporting hyperplane for $D^2/H$]
\label{prop:oa}
Let $(D^{*},H^{*})$ with $H^{*}>0$ be any reference point.  Then for
all $(D,H)$ with $H>0$,
\begin{equation}
  \frac{D^{2}}{H}
  \;\geq\;
  \frac{2D^{*}}{H^{*}} \cdot D
  \;-\;
  \left(\frac{D^{*}}{H^{*}}\right)^{\!2} H.
  \label{eq:oa_cut}
\end{equation}
The bound is tight at $(D,H)=(D^{*},H^{*})$.
\end{proposition}

\begin{proof}
The gradient of $f$ at $(D^{*},H^{*})$ is
$\nabla f = (2D^{*}/H^{*},\,-D^{*2}/H^{*2})$.
Convexity of $f$ gives $f(D,H) \geq f(D^{*},H^{*}) + \nabla
f^{\top}(D-D^{*},H-H^{*})$.  Expanding and simplifying yields
\eqref{eq:oa_cut}. Tightness at the reference point follows by direct
substitution.
\end{proof}

Replacing the SOC constraint \eqref{eq:soc} by an initially empty
set of linear cuts of the form \eqref{eq:oa_cut}—added iteratively at
successive master solutions—converts each master solve into a pure
MILP while preserving the validity of the lower bound.

\subsection{Master Problem}

The master problem at iteration $k$ is:

\begin{subequations}
\label{eq:master}
\begin{align}
  \min \quad
  & \underbrace{g(x,y,z,\eta,\theta)}_{\text{routing objective}}
    + \frac{1}{2K^{2}\mu^{2}(1-\rho)}
      \sum_{r\in\mathcal{R}} \beta_r
    + \text{const}
  \label{eq:master_obj}\\[4pt]
  \text{s.t.} \quad
  & \text{all routing constraints \eqref{eq:assign}--\eqref{eq:visit}},
  \label{eq:master_routing}\\
  & D_r = \sum_{i\in\mathcal{C}} \lambda_i \,\eta_{ir},
    \quad r\in\mathcal{R},
  \label{eq:master_D}\\
  & \beta_r \geq 0,
    \quad r\in\mathcal{R},
  \label{eq:master_beta_lb}\\
  & \beta_r \geq
      \frac{2D_r^{(\ell)}}{H_r^{(\ell)}} D_r
      - \left(\frac{D_r^{(\ell)}}{H_r^{(\ell)}}\right)^{\!2} H_r,
    \quad r\in\mathcal{R},\;\ell=1,\dots,k-1.
  \label{eq:master_cuts}
\end{align}
\end{subequations}

Here $\eta_{ir} = y_{ir} H_r$ is the product of the binary
clinic-assignment variable and the headway (linearised via the
indicator constraints already present in the base formulation);
$(D_r^{(\ell)}, H_r^{(\ell)})$ are the values of $(D_r, H_r)$
extracted from the master solution at iteration $\ell$.  The routing
objective $g$ equals $\frac{1}{|\mathcal{C}|}\sum_{i,r}
(0.5\,\eta_{ir}+\theta_{ir})$ for \textsc{ATAT-Min} and $\max_{i}
\text{TLT}_i$ for \textsc{MTAT-Min}.

Because \eqref{eq:master_cuts} is a finite set of linear constraints,
problem \eqref{eq:master} is a MILP at every iteration.

\subsection{Algorithm}

\begin{algorithm}[t]
\caption{Outer-Approximation L-Shaped Method for the EID-VRP}
\label{alg:lshaped}
\begin{algorithmic}[1]

\State \textbf{Input:} instance data, objective \texttt{obj},
       vehicle mode \texttt{mode}, tolerances $\epsilon_{\text{SOC}}$,
       $\epsilon_{\text{gap}}$, time limit $T_{\max}$.

\State Build master MILP \eqref{eq:master} with
       $\mathcal{L} \leftarrow \emptyset$ (no OA cuts).
\State $\overline{v} \leftarrow +\infty$ \;(best upper bound),\quad
       $\underline{v} \leftarrow -\infty$ \;(best lower bound).

\Repeat

  \State Solve master \eqref{eq:master}; let
         $(\bar{x},\bar{y},\bar{z},\bar{\eta},\bar{\theta},
           \bar{D},\bar{H},\bar{\beta})$ be the incumbent.
  \State $\underline{v} \leftarrow
         \max\!\bigl(\underline{v},\,v^{\text{LB}}_{\text{master}}\bigr)$.
         \Comment{Gurobi's branch-and-bound bound}

  \For{each route $r \in \mathcal{R}$}
    \If{$\bar{H}_r > 0$}
      \State $f_r \leftarrow \bar{D}_r^{2} / \bar{H}_r$
    \Else
      \State $f_r \leftarrow 0$
    \EndIf
  \EndFor

  \State $\varphi_{\text{true}} \leftarrow \sum_r f_r$.
  \State Compute $W_{\text{lab}}$ via \eqref{eq:wlab} using
         $\varphi_{\text{true}}$.
  \State $v_{\text{actual}} \leftarrow
         g(\bar{x},\bar{y},\bar{z},\bar{\eta},\bar{\theta})
         + W_{\text{lab}}$.
  \If{$v_{\text{actual}} < \overline{v}$}
    \State $\overline{v} \leftarrow v_{\text{actual}}$;\quad
           save current routing as best solution.
  \EndIf

  \State $\delta_{\text{SOC}} \leftarrow
         \max_{r}\,\bigl(f_r - \bar{\beta}_r\bigr)$.
  \State $\text{gap} \leftarrow
         (\overline{v} - \underline{v})\,/\,|\overline{v}|$.

  \If{$\delta_{\text{SOC}} \leq \epsilon_{\text{SOC}}$
      \textbf{ and }
      $\text{gap} \leq \epsilon_{\text{gap}}$}
    \State \textbf{break} \Comment{Converged}
  \EndIf

  \For{each $r \in \mathcal{R}$ with $f_r - \bar{\beta}_r > \epsilon_{\text{SOC}}$}
    \State Add to $\mathcal{L}$: \quad
           $\beta_r \;\geq\;
            \dfrac{2\bar{D}_r}{\bar{H}_r}\,D_r
            - \left(\dfrac{\bar{D}_r}{\bar{H}_r}\right)^{\!2} H_r$
    \label{alg:cut}
  \EndFor

\Until{time limit $T_{\max}$ reached}

\State \textbf{Return} best solution found, $\overline{v}$,
       $\underline{v}$, gap, iteration log.

\end{algorithmic}
\end{algorithm}

\subsection{Correctness and Convergence}

\begin{theorem}[Valid lower bound]
At every iteration, the master objective value is a valid lower bound
on the optimal value of the full MIQCP.
\end{theorem}

\begin{proof}
Every OA cut \eqref{eq:oa_cut} is a necessary condition for
$\beta_r \geq D_r^2/H_r$.  Hence the feasible region of the master
contains the projection of the MIQCP feasible set onto $(x,y,z,\beta)$
space at every iteration.  The master therefore minimises over a
superset of feasible solutions, giving a lower bound.
\end{proof}

\begin{theorem}[Finite convergence]
Algorithm~\ref{alg:lshaped} converges in finitely many iterations
when the routing variables $(x,y,z)$ take values in a finite set.
\end{theorem}

\begin{proof}
Given fixed binary $x$ and $y$ and integer $z$, the values of $D_r$
and $H_r$ are fully determined by the linear constraints of the
model.  There are finitely many such combinations.  At each iteration
where the algorithm does not terminate, at least one new OA cut is
added that is violated by the current solution.  Because OA cuts are
never removed, the same $(D_r,H_r)$ combination cannot repeat without
the cut being satisfied, so the algorithm visits each combination at
most once before termination.
\end{proof}

\subsection{Relationship to the L-Shaped Method}

The algorithm follows the spirit of the L-shaped decomposition of
\citet{VanSlyke1969} and its extension to integer programmes by
\citet{Laporte1993}.  In the classical stochastic VRP setting
\citep{Laporte1993,Laporte2002}, the second stage evaluates a
recourse function $Q(y)$ that is convex in the first-stage variables
$y$; optimality cuts $\theta \geq Q(y^*)+ \nabla Q(y^*)^\top(y-y^*)$
are added to the master.  Our setting is deterministic but nonlinear:
the ``recourse'' is the evaluation of $\varphi = \sum_r D_r^2/H_r$,
and the supporting hyperplane of Proposition~\ref{prop:oa} plays the
role of the optimality cut.  The key structural advantage is that the
cut is \emph{linear} in the master variables $D_r$ and $H_r$—both
already present in the routing model—so no auxiliary subproblem needs
to be solved; the cut coefficients are computed in closed form from the
master solution.

\subsection{Implementation Details}

The master model is built once before the loop and modified
in-place by appending OA constraints (line~\ref{alg:cut} of
Algorithm~\ref{alg:lshaped}); Gurobi's incremental modelling interface
avoids rebuilding the branch-and-bound tree from scratch.  Gurobi's
internal bound ($v^{\text{LB}}_{\text{master}}$) from the most recent
solve is used as the lower bound update, which is tighter than the
master's objective value when the solve is terminated by a time limit.

The algorithm is terminated when both criteria hold simultaneously:
(i) the maximum SOC violation
$\delta_{\text{SOC}} = \max_r (f_r - \bar\beta_r) \leq 10^{-4}$,
ensuring that all $\beta_r$ values correctly represent $D_r^2/H_r$;
and (ii) the optimality gap
$(\overline{v}-\underline{v})/|\overline{v}| \leq 10^{-3}$,
ensuring near-global optimality.

For routes with $\bar{H}_r = 0$ (empty routes carrying no clinics),
the term $D_r^2/H_r$ is zero by convention and no cut is added.

\medskip
\noindent
The implementation is available as \texttt{lshaped\_solver.py} and
returns the same structured output as the direct Gurobi solver,
making it a drop-in replacement in the experimental pipeline.

% ── Bibliography entries to add to your .bib file ────────────────────────────
% @article{VanSlyke1969,
%   author  = {Van Slyke, Richard M. and Wets, Roger},
%   title   = {L-shaped linear programs with applications to optimal
%               control and stochastic programming},
%   journal = {SIAM Journal on Applied Mathematics},
%   volume  = {17}, number = {4}, pages = {638--663}, year = {1969}
% }
% @article{Laporte1993,
%   author  = {Laporte, Gilbert and Louveaux, Fran\c{c}ois V.},
%   title   = {The integer {L}-shaped method for stochastic integer
%               programs with complete recourse},
%   journal = {Operations Research Letters},
%   volume  = {13}, number = {3}, pages = {133--142}, year = {1993}
% }
% @article{Laporte2002,
%   author  = {Laporte, Gilbert and Louveaux, Fran\c{c}ois V.
%               and van Hamme, Luc},
%   title   = {An integer {L}-shaped algorithm for the capacitated
%               vehicle routing problem with stochastic demands},
%   journal = {Operations Research},
%   volume  = {50}, number = {3}, pages = {415--423}, year = {2002}
% }
